\subsection{包含关系与引用的转换}

对 \texttt{EReference} 的处理取决于其功能。\texttt{input\_value} 注解生成
水平连接；带有 \texttt{containment=true} 的引用生成结构输入；其余引用则
生成动态选择器。代码清单 \ref{lst:conversion-containment-ecore} 仅保留
\texttt{App.pages} 所使用的分支，省略了不会改变这一结构创建过程的
其他分支。

在该代码片段中，\texttt{eRef} 是输入的 \texttt{EReference}，
\texttt{refAnn} 是其可选的 Blockly \texttt{EAnnotation}。\texttt{bt} 是所属
类的 \texttt{BlockTypeSpec}，\texttt{sis} 是输出的
\texttt{Statement}\allowbreak\texttt{InputSpec}。\texttt{check} 保存连接
所接受的类型；\texttt{hasContainments} 记录该块包含子元素。数据流为
\texttt{eRef} $\rightarrow$ \texttt{sis} $\rightarrow$ \texttt{bt}。

\begin{lstlisting}[style=tfmjavaecore, captionpos=t, caption={Ecore 包含关系到结构输入的转换}, label={lst:conversion-containment-ecore}]
if (eRef.isContainment()) {
    hasContainments = true;
    StatementInputSpec sis = new StatementInputSpec();
    sis.setName(eRef.getName());
    String check = refAnn != null
        ? refAnn.getDetails().get("check") : null;
    // 在 sis 中配置显式类型，或使用 eRef 的 EClass。
    sis.setCheckType(check != null && !check.isBlank()
        ? check
        : eRef.getEReferenceType().getName());
    sis.setLowerBound(eRef.getLowerBound());
    // 通用模型以零表示开放上界。
    sis.setUpperBound(eRef.getUpperBound() == -1
        ? 0 : eRef.getUpperBound());
    bt.getStatementInputs().add(sis);
}
\end{lstlisting}

因此，对于 \texttt{App.pages}，适配器会创建一个名为
\texttt{pages} 的 \texttt{Statement}\allowbreak\texttt{InputSpec}，其接受类型为
\texttt{Page}，下界为一，上界为五。\texttt{Page} 块具有与该输入兼容的
连接。同一转换还把 \texttt{Page.components} 映射为接受
\texttt{Component} 的语句输入；\texttt{Label}、\texttt{Button} 和
\texttt{Image} 通过继承满足该类型。基数还用于生成验证规则，因此数量超出
范围时编辑器会发出警告。

\subsection{Ecore 映射总结}

\begin{table}[H]
  \centering
  \small
  \begin{tabular}{|L{0.30\textwidth}|L{0.30\textwidth}|L{0.26\textwidth}|}
    \hline
    输入 Ecore 元素 & 目标 \texttt{EditorSpec} 元素 & 在 Blockly 中的用途 \\
    \hline
    \texttt{EPackage} & 包含名称、\texttt{nsURI} 和全局选项的
    \texttt{EditorSpec} & 编辑器标识与 XMI 导出。 \\
    具体 \texttt{EClass} & \texttt{BlockTypeSpec} & 工具箱中可用的块类型。 \\
    抽象 \texttt{EClass} & 抽象 \texttt{BlockTypeSpec} & 用于连接但不可实例化的
    类型。 \\
    \texttt{EAttribute} & \texttt{FieldSpec} & 文本、数字、布尔、颜色、角度或
    下拉字段。 \\
    \texttt{EReference containment=true} & \texttt{StatementInputSpec} & 用于
    子块的垂直输入。 \\
    \texttt{EReference containment=false} & \texttt{ReferenceFieldSpec} & 指向
    现有元素的动态下拉框。 \\
    \texttt{lowerBound}/\texttt{upperBound} & \texttt{ValidationSpec} 与输入
    限制 & 必填性与基数警告。 \\
    \texttt{EAnnotation} & 块、界面、验证或代码元数据 &
    颜色、类别、标签、工具提示、控件和模板。 \\
    \hline
  \end{tabular}
  \caption{从 Ecore 到 \texttt{EditorSpec} 的主要转换}
  \label{tab:transformacion-ecore-editor-spec}
\end{table}

\subsection{连接与继承的推断}

连接推断同样在适配器中完成。包含其他元素且本身不被其他元素包含的
类可以转换为没有垂直连接的根块。继承自超类的类可以生成类型化连接。
标记为输出块的类会生成值连接器。这种推断很重要，因为用户无需为每个
类手工编写 Blockly 连接规则。在第一个示例中，\texttt{App} 被识别为
根块，\texttt{Page} 获得与 \texttt{App.pages} 兼容的连接，
\texttt{Label}、\texttt{Button} 和 \texttt{Image} 则获得与
\texttt{Component} 兼容的连接。

\section{来自文本路线 \texttt{.m2b} 的适配器}
\label{sec:implementacion-domainmodel-adapter}

文本路线的适配器同样由 Java 手工实现。
\texttt{DomainModelAdapter.fromDomainModel} 接收 Xtext 生成的 EMF 根对象，
并依次执行 \texttt{convertCategories}、\texttt{convertClasses}、
\texttt{convertValidations} 和 \texttt{convertConstraints}。随后，它生成
由基数、必填性和唯一性推断出的验证规则，最后通过
\texttt{convertWorkspace}\allowbreak\texttt{Options} 复制工作区选项。
所得 Java 对象随后返回到代码清单 \ref{lst:entrada-adaptadores-java} 中的
\texttt{toEditorSpec} 方法，并在那里转换为与 Ecore 路线相同的通用 EMF
模型。

该映射比 Ecore 更直接，因为语法已经区分了 Model2Blockly 所需的特定
概念。属性转换为字段，包含关系转换为结构输入，引用转换为动态下拉框，
值输入转换为水平连接。顺序约束转换为验证规则，标记为必填的字段则生成
必填性检查。

每个类的转换通过遍历 Xtext 生成的 \texttt{Feature} 列表来实现。
代码清单 \ref{lst:despacho-domainmodel-adapter} 展示了核心分派逻辑。
每个具体类型都会交给专用方法处理，输入名称则加入
\texttt{orderedInputNames}；这样，在构造块时可以保留 \texttt{.m2b}
文件中的书写顺序。

\begin{lstlisting}[style=tfmjava, captionpos=t, caption={\texttt{DomainModelAdapter} 中 DSL 元素的分派}, label={lst:despacho-domainmodel-adapter}]
for (Feature feature : cls.getFeatures()) {
    if (feature instanceof Attribute) {
        FieldSpec field = convertAttribute((Attribute) feature);
        bt.getFields().add(field);
        bt.getOrderedInputNames().add(field.getName());
    } else if (feature instanceof Containment) {
        StatementInputSpec input =
            convertContainment((Containment) feature);
        bt.getStatementInputs().add(input);
        bt.getOrderedInputNames().add(input.getName());
    } else if (feature instanceof Reference) {
        ReferenceFieldSpec reference =
            convertReference((Reference) feature);
        bt.getReferenceFields().add(reference);
        bt.getOrderedInputNames().add(reference.getName());
    } else if (feature instanceof ValueInput) {
        ValueInputSpec input = convertValueInput((ValueInput) feature);
        bt.getValueInputs().add(input);
        bt.getOrderedInputNames().add(input.getName());
    }
}
\end{lstlisting}

\begin{table}[H]
  \centering
  \small
  \begin{tabular}{|L{0.30\textwidth}|L{0.30\textwidth}|L{0.26\textwidth}|}
    \hline
    DSL 元素 & 目标 \texttt{EditorSpec} 元素 & 在 Blockly 中的用途 \\
    \hline
    \texttt{domain} & \texttt{EditorSpec} & 编辑器名称与生成选项。 \\
    \texttt{category} & \texttt{CategorySpec} & 工具箱类别与默认颜色。 \\
    \texttt{class} & \texttt{BlockTypeSpec} & 块类型。 \\
    \texttt{attribute} & \texttt{FieldSpec} & 块的可视字段。 \\
    \texttt{contains [m..n]} & \texttt{StatementInputSpec} 与
    \texttt{ValidationSpec} & 子元素输入与基数检查。 \\
    \texttt{reference} & \texttt{ReferenceFieldSpec} & 动态下拉框。 \\
    \texttt{value} & \texttt{ValueInputSpec} & 值块的水平输入。 \\
    \texttt{must follow} & \texttt{ValidationSpec} & 局部顺序规则。 \\
    \hline
  \end{tabular}
  \caption{将文本 DSL \texttt{.m2b} 规范化为 \texttt{EditorSpec}}
  \label{tab:transformacion-dsl-editor-spec}
\end{table}

实现还利用继承信息来确定连接形态。标记为输出块的类会转换为值块；
具有超类的类会生成类型化连接；具有包含关系且没有超类的类可以作为
根块；普通类则生成为可堆叠块。通过这种方式，文本语法中声明的语义会
转换为 Blockly 可视规则，领域作者无需了解该库的底层 API。

系统使用 Xtext 生成的 Java 类型（\texttt{ClassDef}、
\texttt{Attribute}、\texttt{Containment}、\texttt{Reference} 和
\texttt{ValueInput}）读取主要结构。只有部分界面、代码和工作区选项
通过反射获取：\texttt{eClass()} 使用 \texttt{getEStructuralFeature}
定位特征，\texttt{eGet} 则获取其值。这一决定使语法中新加入的元数据
可以在重新生成全部访问器之前被读取，同时仍能在确定块结构的遍历过程
中保留静态检查。

\section{HTML/JavaScript 生成器}
\label{sec:implementacion-generador-html}

HTML/JavaScript 输出生成器将中间模型转换为在浏览器中打开编辑器所需
的产物：块、工具箱、导出、验证、编辑页面、独立页面、验证工作区、
示例模型和生成报告。该输出使用 Blockly 定义块、注入工作区、管理
工具箱、序列化并生成代码 \cite{blocklyDocs}。

代码生成主要使用 Xtend 实现。主类
\texttt{Blockly\allowbreak{}Code\allowbreak{}Generator.xtend} 接收一个
\texttt{EditorSpec}，并生成从文件名到文本内容的有序映射。代码清单
\ref{lst:entrada-generador-xtend} 展示了入口方法：每个
\texttt{generate*} 函数返回一个字符串，该映射则保存该字符串与其目标
文件之间的关系。

\begin{lstlisting}[style=tfmxtend, captionpos=t, caption={\texttt{BlocklyCodeGenerator.xtend} 的生成入口}, label={lst:entrada-generador-xtend}]
def Map<String, String> generate(EditorSpec spec) {
    val files = new LinkedHashMap<String, String>()
    // 获取基础名称；若未定义，则使用 domain。
    val base = spec.domainName ?: "domain"
    files.put(base + "_blocks.js", generateBlocksJs(spec))
    files.put(base + "_toolbox.js", generateToolboxJs(spec))
    files.put(base + "_generators.js", generateGeneratorsJs(spec))
    files.put(base + "_validations.js", generateValidationsJs(spec))
    files.put(base + "_editor.html", generateEditorHtml(spec))
    files.put(base + "_standalone.html", generateStandaloneHtml(spec))
    files.put("validation_workspace.html",
        ValidationWorkspaceHtmlGenerator.generate(spec))
    files.put("validation_blocks.json",
        ValidationBlockModelGenerator.generate(spec))
    files.put("validation_runtime.js",
        ValidationRuntimeGenerator.generateClassicJs(spec))
    files.put("sample_model.json", SampleModelGenerator.generate(spec))
    return files
}
\end{lstlisting}

选择 Xtend 作为生成机制，是因为它与 Xtext 集成良好，并可在 Eclipse
项目中编写易读的文本模板。本项目没有使用 Acceleo 或外部模板引擎。
示例模型、验证运行环境和规则可视化视图等辅助部分由单独的 Java 类
实现，以保持各自职责边界清晰。

该映射不会直接写入磁盘。\texttt{GenerateBlocklyEditorHandler} 在其
路径前添加 \texttt{html/}，并加入三个在 Xtend 之外生成的结果：中间
XMI、\texttt{README.md} 和生成报告。最后，\texttt{writeFiles} 创建
工作区中的 \texttt{IFile} 并刷新项目。这种分离使同一生成器可由
独立入口复用，而无需依赖 Eclipse API。

报告由 \texttt{GenerationReport}\allowbreak\texttt{HtmlRenderer}
构建。该类接收源文件、\texttt{EditorSpec} 和产物清单，并生成
\texttt{generation\_}\allowbreak\texttt{report.html}。
\texttt{README.md} 则在处理器或相应独立入口中组装，因为它包含针对
执行方式定制的说明。这两个结果都会在写入之前加入映射，但不属于
\texttt{BlocklyCodeGenerator} 的直接输出。

块生成从 \texttt{generateBlocks}\allowbreak\texttt{Array} 开始；该方法
筛选具体类型并调用
\texttt{blockTypeTo}\allowbreak\texttt{BlockJson}。后者创建字段、引用和输入的映射，并遍历
\texttt{orderedInput}\allowbreak\texttt{Names}，以保留适配器确定的顺序。每个元素分别交由
\texttt{fieldToArg}、\texttt{referenceToArg}、
\texttt{valueInputToArg} 或
\texttt{statementInput}\allowbreak\texttt{ToArg} 处理。最后，对
\texttt{Connection}\allowbreak\texttt{Type} 执行的
\texttt{switch} 会生成自由类型或类型化的
\texttt{previous}\allowbreak\texttt{Statement}、
\texttt{nextStatement} 或 \texttt{output}。抽象类保留在模型中，以便
解析继承和连接，但不会作为可实例化块加入输出。

工具箱生成过程会产出用户可用的块集合。若没有类别，生成器可以输出
简单工具箱；若存在类别，则生成带类别的工具箱，并在已定义子类别时将
其一并纳入。还会加入 Blockly 的内置类别，如逻辑、循环、数学、变量
和函数，使各领域可在必要时将领域块与通用构造结合使用。
